\chapter{Cross-Dataset Validation on UNSW-NB15}
\label{ch:crossdataset}

The CICIDS2017 results so far are internally consistent: random forest
is the most temporally robust binary detector, and the binary task
survives temporal shift while the multi-class task does not. The
question this chapter addresses is whether those conclusions transfer
to a different dataset --- specifically, UNSW-NB15, generated by a
different team using a different testbed and a substantially different
attack catalogue.

Table~\ref{tab:unsw} reports UNSW-NB15 macro F1 for the multi-class task
and F1 for the binary task using the official train/test split
(175{,}341 test rows). The model ranking flips in the multi-class
setting: LightGBM leads at 0.548, XGBoost trails at 0.522, random forest
is third at 0.479, and logistic regression is last at 0.404. The binary
setting is much more uniform, with all three tree models clustered
around F1 $\approx 0.92$ and ROC-AUC $\approx 0.986$.

\begin{table}[H]
    \centering
    \caption{UNSW-NB15 macro F1 (multi-class) and F1 (binary) using the
    official train/test split.}
    \label{tab:unsw}
    \small
    \begin{tabular}{lccc}
        \toprule
        \textbf{Model} & \textbf{Multi-class F1} & \textbf{Binary F1} & \textbf{Binary ROC-AUC} \\
        \midrule
        LogReg & 0.404 & 0.888 & 0.976 \\
        Random Forest & 0.479 & 0.921 & 0.986 \\
        XGBoost & 0.522 & 0.919 & 0.986 \\
        LightGBM & \textbf{0.548} & 0.921 & 0.986 \\
        \bottomrule
    \end{tabular}
\end{table}

\begin{figure}[H]
    \centering
    \includegraphics[width=0.95\textwidth]{cross_dataset_comparison.png}
    \caption{Model performance side-by-side on CICIDS2017 versus
    UNSW-NB15. The model ranking flips between datasets in the
    multi-class setting --- XGBoost wins on CICIDS multi-class but
    LightGBM wins on UNSW.}
    \label{fig:cross-comparison}
\end{figure}

The rank reversal between datasets is the most important cross-dataset
finding. No single tree ensemble dominates universally, and the headline
choice of ``best model'' depends on which dataset the benchmark is run
on. Practitioners who deploy a model based on a single-dataset result
should be aware that the same model may perform poorly on traffic from a
different testbed, and that validation on at least two independent
datasets is a minimum quality bar before claiming a universal winner.

The mechanism behind the rank reversal is feature-distribution
mismatch. On CICIDS, the top-ranked features for random forest are
inter-arrival-time statistics, because CICIDS attacks differ most from
benign traffic in their timing structure. On UNSW, the top-ranked
features are protocol and service categorical features
(\texttt{service\_dns}, \texttt{state\_INT}, \texttt{ct\_dst\_sport\_ltm}),
because UNSW attacks differ most from benign traffic in which services
they target and how connection states evolve. A model that best exploits
the dominant feature family on one dataset is not the best on the other.
